% !TeX root = ../thuthesis-example.tex

\chapter{结论与讨论}

\section{结论}
\label{sec:conclusion}

本研究的结论以第6.9节结果表中的核心指标为依据：Deletion-AUC、Insertion-AUC、NDCG@Ke-AUC、相对随机基线增益（minus-random）、POS@10、运行耗时以及多随机种子下的统计检验（Wilcoxon 与 Cliff's $\delta$）。据此可归纳如下。

其一，当前结果支持如下\textbf{趋势性发现}：SIF/SIF+ 在计算开销与归因效果之间具有较优折中，能够在统一评测预算下提供可解释且可部署的样本级价值估计。

其二，考虑到当前样本规模等实验条件下部分比较尚未达到校正后显著性门槛，本文不将趋势性差异外推描述为确定性优越结论。

其三，本文把 VAE 与 NCF 作为后续完善验证模型，并在不变的统计口径下持续追加结果；后续版本仅在“方向一致 + 校正后显著 + 效应量充分”的条件下升级为强结论。

综上，本文当前提交版本完成了“理论框架 + 实证趋势 + 复现流程”的闭环，并为后续实证扩展提供了可审计、可持续更新的研究路径。

\section{应用前景与讨论}
\label{sec:discussion}

\subsection{方法局限性分析}

SIF 与 SIF+ 的理论推导均依赖模型参数收敛至稳定点 $\theta^*$ 的假设（假设 A2）。在此假设下，验证梯度 $\nabla_\theta L_{\mathrm{val}}(\theta^*)$ 与样本梯度的内积能够稳定反映样本对泛化目标的边际贡献。然而，工业场景中出于防止过拟合或节省算力的考量，模型往往采用\textbf{早停策略}（early stopping），在损失尚未充分收敛时即终止训练。此时参数 $\theta_t$ 距离稳定点存在系统性偏差，导致以下问题：

\begin{itemize}
    \item \textbf{梯度方向不稳定：}早停点附近的验证梯度仍处于快速变化阶段，单次快照的内积估计方差较大，SIF 分数的跨样本排序可靠性下降。
    \item \textbf{曲率近似失效：}SIF+ 的对角 Hessian 近似在收敛邻域内具有理论保证，但在早停点处 Hessian 特征值分布可能尚未稳定，对角近似的误差界不再成立。
    \item \textbf{阶段分解偏移：}DVF 的阶段分解以梯度范数阈值 $\epsilon$ 划定表示学习与对齐阶段的边界，早停可能导致对齐阶段贡献被截断，低估后期训练样本的价值。
\end{itemize}

针对上述局限，一种可行的缓解策略是在早停点之后继续执行若干步"价值评估专用微调"（valuation fine-tuning），以较小学习率使模型在验证集上进一步收敛，再计算 SIF/SIF+ 分数；或采用多检查点平均策略，对早停前后若干检查点的 SIF 分数取加权均值，以降低单点估计的方差。

\subsection{商业应用价值}

本文框架的应用需严格遵循相关法律法规：
\begin{itemize}
    \item 《中华人民共和国个人信息保护法》第四条、第六条：原始行为日志（含user id）属于个人信息，但特征级embedding、统计特征、偏好向量经匿名化处理后不属于个人信息；本文不对完整用户画像整体定价，符合最小必要原则。
    \item 《中华人民共和国数据安全法》第二十七条：本文方案不对原始明细数据做外部流通，仅输出价值仪表盘，保障数据安全。
    \item 《互联网信息服务算法推荐管理规定》第十六条：归因过程使推荐算法的决策权重可追溯，增强了算法透明度。
\end{itemize}

在合规前提下，该框架可在以下三个维度释放商业价值：

\begin{itemize}
    \item \textbf{数据采购决策：}通过 SIF+ 分数量化不同来源数据对推荐效果的边际贡献，为数据采购谈判提供客观定价依据，避免"按量付费"模式下的价值错配。
    \item \textbf{数据资产仪表盘：}周期性输出样本级、实体级与数据集级价值报告，回答"本次促销增长主要归功于哪些数据"、"每TB用户行为数据的季度边际收益是多少"等具体管理问题，支撑数据要素的合规高效流通。
    \item \textbf{数据清洗与治理：}将持续低价值或负价值样本（SIF+ 分数显著为负）标记为噪声候选，辅助数据质量团队制定清洗优先级，在不损失模型效果的前提下降低存储与计算成本。
\end{itemize}

\subsection{未来工作展望}
本文实验部分正在持续完善中，后续版本将补充以下内容：
\begin{itemize}
    \item \textbf{模型多样性：}除 MLP 外，增加 VAE 与 NCF 作为验证模型，以测试方法在不同架构下的适用性与稳定性。
    \item \textbf{数据规模扩展：}在更大规模数据集（如 MovieLens 20M）上验证方法的可扩展性与效果趋势。
    \item \textbf{早停场景验证：}针对早停策略下的估值稳定性问题，设计专门的实验验证上述局限性分析中的缓解方案。
    \item \textbf{商业应用A/B测试：}在实际推荐系统中部署 SIF+ 作为数据价值监控工具，评估其对数据采购决策与清洗优先级制定的实际影响。
\end{itemize}

此外，未来工作还将探索更高阶的曲率近似方法（如 K-FAC）在 SIF+ 中的应用，以及将 DVF 与数据 Shapley 值的理论联系进一步形式化，以深化对数据价值评估方法之间关系的理解。积分梯度等公理化归因方法\cite{sundararajan2017axiomatic}与本文框架在归因公理层面存在潜在联系，亦是值得探索的理论延伸方向。本文实验代码与数据处理流程已开源\cite{repository}，以支持后续复现与扩展研究。

